\scanpage{317}
\mspage{233}
% archive sheet number: 316

\unsure{n'a} pas forcément $f(w)=f(v)u'$, on a là une double flèche \deleted{\unsure{dans I}} $f(i_0)\unsure{=|\rightrightarrows}f(i_2)$ dans $I'$. Mais

\deleted{$\unsure{\zeta|\xi}$ \inserted{$ {}_{f(i_0)}\backslash I$} étant filtrante, cette double flèche peut être égalisée, les deux objets $\zeta'_2=(i_2,\deleted{\unsure{\zeta}}f(w)),\zeta''_2=(i_2,$}

par la condition $\mathrm{F}2$ sur $f$, il existe $i_2\xrightarrow{\alpha}\overline{i}_2$ dans $I$ telle que $f(\alpha)$ égalise ladite double flèche, donc $f(\underbrace{\alpha w}_{\overline{w}})=f(\underbrace{\alpha v}_{\overline{v}})u'$, et quitte à remplacer $i_2$ par $\overline{i}_2$, et \deleted{\unsure{s}} $w,v$ par $\overline{w}=\alpha w$, $\overline{v}=\alpha\unsure{w|v}$, on \unsure{gagne}.

Reste à prouver (pour la partie 1\textsuperscript{o}), que si $f$ est \uline{fidèle}, alors $c)\Rightarrow\deleted{\unsure{a}}b)$. Il reste donc à prouver que \unsure{b)|c)} implique que \deleted{\unsure{pour}} pour $i'$ dans $I'$, $ {}_{i'}\backslash I$ est filtrante. On sait déjà par $\mathrm{F}1$ que $ {}_{i'}\backslash I$ est non vide, il reste à prouver les conditions $\mathrm{PF}1'$, $\mathrm{PF}2$ pour $ {}_{i'}\backslash I$. Comme $I'$ \unsure{filtrant|filtrante}, déjà on peut compléter

\diagram{\begin{tikzcd}[ampersand replacement=\&,column sep=small,row sep=tiny]
\&f(i_1)\arrow[dr,dashed,no head]\&\\
i'\arrow[ur,no head]\arrow[dr,no head]\&\&\unsure{j''|j'}\\
\&f(i_2)\arrow[ur,dashed,no head]\&
\end{tikzcd}}

\diagram{\begin{tikzcd}[ampersand replacement=\&,row sep=tiny,column sep=tiny]
\&i_1\arrow[rr,dashed,no head,"v_1?"]\arrow[dd,no head]\&\&i_3\arrow[r]\arrow[dd,no head]\&\overline{i}_3\arrow[dd,no head]\\
\&i_2\arrow[urr,no head,"v_2?"']\arrow[dd,no head]\&\&\&\\
\&f(i_1)\arrow[rr,no head,"v'_1"]\&\&f(i_3)\arrow[r,no head]\&f(\unsure{i'_3|\overline{i}_3})\\
i'\arrow[ur,no head,"u'_1"]\arrow[r,no head,"u'_2"']\&f(i_2)\arrow[urr,"v'_2"']\&\&\&
\end{tikzcd}}

, et \deleted{\unsure{comme}} par $\mathrm{F}1$ OPS $j'$ de la forme $f(i_3)$. Appliquant $\mathrm{F}3$ à $i_1,i_3,v'$, OPS (quitte à remplacer $i_3$ par $\overline{i}_3$)

